\documentclass{article}
\usepackage[UTF8]{ctex}
\usepackage{amsmath,amssymb,amsfonts}
\usepackage{enumitem}
\usepackage{geometry}
\geometry{a4paper, margin=1in}
% === 错题本排版优化（不参与题目标识） ===
\setlength{\parindent}{0pt}
\setlength{\parskip}{0.35em plus 0.10em minus 0.08em}
\setlength{\emergencystretch}{3em}
\setlist{itemsep=0.25em, topsep=0.40em, parsep=0pt}
\allowdisplaybreaks[2]
% === 错题本排版优化结束 ===
\begin{document}
\section*{西北工业大学2025年高等代数真题}
\section*{172．西北工业大学2025年高等代数试卷}

\subsection{题目1}
1．（15 分）计算多项式函数一阶导数 \(f^{\prime}(x)\) ，其中

\[
f(x)=\left|\begin{array}{cccc}
a_{11}+x & a_{12}+x & \cdots & a_{1 n}+x \\
a_{21}+x & a_{22}+x & \cdots & a_{2 n}+x \\
\vdots & \vdots & & \vdots \\
a_{n 1}+x & a_{n 2}+x & \cdots & a_{n n}+x
\end{array}\right|, n \geq 2 .
\]

\subsection{题目2}
2．（15 分）已知

\[
A=\left(\begin{array}{ccc}
1 & a & b \\
-a & 1 & c \\
-b & -c & 1
\end{array}\right)
\]

是 3 阶实方阵。证明：有关 \(X\) 的矩阵方程 \(A^{2} X-X A-X=O\) 只有零解。

\subsection{题目3}
3．（15 分）齐次线性方程组

\[
\left\{\begin{array}{l}
x_{1}+2 x_{2}+3 x_{3}=0 \\
2 x_{1}+3 x_{2}+5 x_{3}=0 \\
x_{1}+x_{2}+a x_{3}=0
\end{array}\right.
\]

和

\[
\left\{\begin{array}{l}
x_{1}+b x_{2}+c x_{3}=0 \\
2 x_{1}+b^{2} x_{2}+(c+1) x_{3}=0
\end{array}\right.
\]

同解，求 \(a, b, c\) 。

\subsection{题目4}
4．（15 分）\(V\) 是有理数域 \(\mathbb{Q}\) 上的 3 维线性空间， \(\mathcal{A}\) 为 \(V\) 上的线性变换，满足

\[
\mathcal{A}(\alpha)=\beta, \mathcal{A}(\beta)=\gamma, \mathcal{A}(\gamma)=\alpha+\beta .
\]

证明：若 \(\alpha \neq 0\) ，则 \(\alpha, \beta, \gamma\) 线性无关．

\subsection{题目5}
5．（15 分）求 6 阶实方阵

\[
A=\left(\begin{array}{cccccc}
a & -b & & & & \\
-b & a & 1 & & & \\
& & a & -b & & \\
& & -b & a & 1 & \\
& & & & a & -b \\
& & & & -b & a
\end{array}\right)
\]

的 Jordan 标准形．

\subsection{题目6}
6．（15 分）设 \(A=\left(a_{i j}\right)\) 是 3 阶实对称方阵，\(A^{*}\) 是 \(A\) 的伴随矩阵，\(E\) 是 3 阶单位矩阵。

\[
f\left(x_{1}, x_{2}, x_{3}\right)=\left|\begin{array}{cccc}
0 & x_{1} & x_{2} & x_{3} \\
-x_{1} & a_{11} & a_{12} & a_{13} \\
-x_{2} & a_{21} & a_{22} & a_{23} \\
-x_{3} & a_{31} & a_{32} & a_{33}
\end{array}\right| .
\]

若行列式 \(\operatorname{det} A=-12\) ，迹 \(\operatorname{tr}(A)=1,(1,0,-2)^{\mathrm{T}}\) 是线性方程组 \(\left(A^{*}-4 E\right) X=0\) 的解．求正交变换\(X=Q Y\) ，将 \(f\left(x_{1}, x_{2}, x_{3}\right)\) 化为标准形、

\subsection{题目7}
7．（15 分）Euclid 空间 \(V\) 上线性变换 \(\mathcal{A}\) 保持基 \(\alpha_{1}, \alpha_{2}, \cdots, \alpha_{n}\) 向量长度不变，即

\[
\left(\mathcal{A} \alpha_{i}, \mathcal{A} \alpha_{i}\right)=\left(\alpha_{i}, \alpha_{i}\right)
\]

问： \(\mathcal{A}\) 是否一定是正交变换？如果是，证明之；如果不是，举出反例．

\subsection{题目8}
8．（15 分）在 \(F^{4}\) 中，\(V_{1}\) 是由

\[
\alpha_{1}=(1,0,-1,0), \alpha_{2}=(0,0,1,-1), \alpha_{3}=(1,1,0,0)
\]

生成的子空间，\(V_{2}\) 是由

\[
\beta_{1}=(1,2,-1,2), \beta_{2}=(0,1,-1,0), \beta_{3}=(0,2,1,-1)
\]

生成的子空间，分别求出 \(V_{1}+V_{2}, V_{1} \cap V_{2}\) 的维数和一组基。

\subsection{题目9}
9．（15 分）设 \(f(x)\) 是实系数二次多项式，且有不同的实根 \(x_{1}, x_{2}, \mathcal{A}\) 是实线性空间 \(V\) 上的线性变换，\(\mathcal{A} \neq x_{1} \mathcal{E}, \mathcal{A} \neq x_{2} \mathcal{E}\)（ \(\mathcal{E}\) 为单位变换），且满足＇\(f(\mathcal{A})=\mathcal{O}\) 。证明：

\begin{itemize}
  \item[（1）] \(x_{1}, x_{2}\) 是 \(\mathcal{A}\) 的特征值．
  \item[（2）] \(V=V_{1} \oplus V_{2}, V_{i}\) 是 \(\mathcal{A}\) 的属于 \(x_{i}\) 的特征子空间．
\end{itemize}

\subsection{题目10}
10．（15 分）设 \(A=\left(a_{i j}\right), B=\left(b_{i j}\right)\) 为 \(n\) 阶方阵，定义 \(A \circ B=\left(a_{i j} b_{i j}\right)\) 。证明：

\begin{itemize}
  \item[（1）] 若 \(A, B\) 都是半正定的，则 \(A \circ B\) 也是半正定的．
  \item[（2）] 若 \(A, B\) 都是正定的，则 \(A \circ B\) 也是正定的．
\end{itemize}
\end{document}
